% GENERATED FILE. Edit canonical lessons, then run npm run generate:books.
% canonical-source-hash: fnv1a64:7beba3fcd2a83e05
% canonical-lessons: ES-C636-mentir, ES-C636-mentira, ES-C636-enganar, ES-C636-honesto, ES-C636-merecer

\chapter{Decir la verdad}
\label{ch:es-decir-la-verdad}
\hlchaptermodality{636}

\begin{chapteropening}
\textbf{By the end of this chapter:} \emph{I can talk about lies, honesty and what someone deserves.}

Complete the last lesson of chapter 636: I can talk about lies, honesty and what someone deserves.
\end{chapteropening}

\section[mentir]{mentir — to lie, to tell lies}

\label{lesson:ES-C636-mentir}



\begin{quote}
Before the new one: say the Spanish for the consistency, then the Spanish for willing.
\end{quote}

\subsection*{You'll want to know: mentir}
\textbf{mentir} — \textquotedblleft{}to lie, to tell lies\textquotedblright{}.

\begin{quote}
No me mientas.
\end{quote}

\begin{quote}
\emph{Do not lie to me.}
\end{quote}

Like \emph{sentir}: \emph{miento, mintió}.

\begin{cousinweb}
\emph{mentir} is inherited from Latin \emph{mentiri} 'to lie, to invent', usually linked with \emph{mens} 'mind' (EtymDB).
\end{cousinweb}

\begin{grammarlens}[title={What you've built}]
One more word for telling, explaining and describing at B1.
\end{grammarlens}

\subsection*{Your turn}
\begin{itemize}
\raggedright
  \item \emph{Say it:} \emph{mentir}
  \item \emph{Say it:} \emph{mentir}, once more
  \item \emph{Say it:} a sentence of your own with \emph{mentir}
  \item \emph{Recall:} say the Spanish for the consistency, then the Spanish for willing, then say \emph{mentir} again
\end{itemize}

\begin{tcolorbox}[breakable,colback=teal!4,colframe=teal!35!black,title={Before you move on}]
What does mentir mean? (\textquotedblleft{}To lie, to tell lies\textquotedblright{}.) Say it once more.
\end{tcolorbox}
\section[la mentira]{la mentira — the lie}

\label{lesson:ES-C636-mentira}



\begin{quote}
Before the new one: say the Spanish for willing, then the Spanish for to lie.
\end{quote}

\subsection*{You'll want to know: la mentira}
\textbf{la mentira} — \textquotedblleft{}the lie\textquotedblright{}.

\begin{quote}
Eso es mentira.
\end{quote}

\begin{quote}
\emph{That is a lie.}
\end{quote}

\emph{Parece mentira} — it is hard to believe.

\begin{cousinweb}
\emph{mentira} is built on \emph{mentir}, inherited from Latin \emph{mentiri} 'to lie' (EtymDB).
\end{cousinweb}

\begin{grammarlens}[title={What you've built}]
One more word for telling, explaining and describing at B1.
\end{grammarlens}

\subsection*{Your turn}
\begin{itemize}
\raggedright
  \item \emph{Say it:} \emph{la mentira}
  \item \emph{Say it:} \emph{la mentira}, once more
  \item \emph{Say it:} a sentence of your own with \emph{la mentira}
  \item \emph{Recall:} say the Spanish for willing, then the Spanish for to lie, then say \emph{la mentira} again
\end{itemize}

\begin{tcolorbox}[breakable,colback=teal!4,colframe=teal!35!black,title={Before you move on}]
What does la mentira mean? (\textquotedblleft{}The lie\textquotedblright{}.) Say it once more.
\end{tcolorbox}
\section[engañar]{engañar — to deceive, to cheat}

\label{lesson:ES-C636-enganar}



\begin{quote}
Before the new one: say the Spanish for to lie, then the Spanish for the lie.
\end{quote}

\subsection*{You'll want to know: engañar}
\textbf{engañar} — \textquotedblleft{}to deceive, to cheat\textquotedblright{}.

\begin{quote}
El vendedor me engañó.
\end{quote}

\begin{quote}
\emph{The seller cheated me.}
\end{quote}

\emph{Engañar a} someone. It can also mean to be unfaithful.

\begin{cousinweb}
\emph{engañar} is usually traced to spoken Latin \emph{ingannare} 'to mock, to trick'; EtymDB has no parent for it, and the origin is not settled.
\end{cousinweb}

\begin{grammarlens}[title={What you've built}]
One more word for telling, explaining and describing at B1.
\end{grammarlens}

\subsection*{Your turn}
\begin{itemize}
\raggedright
  \item \emph{Say it:} \emph{engañar}
  \item \emph{Say it:} \emph{engañar}, once more
  \item \emph{Say it:} a sentence of your own with \emph{engañar}
  \item \emph{Recall:} say the Spanish for to lie, then the Spanish for the lie, then say \emph{engañar} again
\end{itemize}

\begin{tcolorbox}[breakable,colback=teal!4,colframe=teal!35!black,title={Before you move on}]
What does engañar mean? (\textquotedblleft{}To deceive, to cheat\textquotedblright{}.) Say it once more.
\end{tcolorbox}
\section[honesto]{honesto — honest}

\label{lesson:ES-C636-honesto}



\begin{quote}
Before the new one: say the Spanish for the lie, then the Spanish for to deceive.
\end{quote}

\subsection*{You'll want to know: honesto}
\textbf{honesto} — \textquotedblleft{}honest\textquotedblright{}.

\begin{quote}
Es una persona honesta.
\end{quote}

\begin{quote}
\emph{She is an honest person.}
\end{quote}

In Spain \emph{honrado} is also common for honest.

\begin{cousinweb}
\emph{honesto} is a learned borrowing from Latin \emph{honestus} 'honourable', built on \emph{honor} 'honour' (EtymDB).
\end{cousinweb}

\begin{grammarlens}[title={What you've built}]
One more word for telling, explaining and describing at B1.
\end{grammarlens}

\subsection*{Your turn}
\begin{itemize}
\raggedright
  \item \emph{Say it:} \emph{honesto}
  \item \emph{Say it:} \emph{honesto}, once more
  \item \emph{Say it:} a sentence of your own with \emph{honesto}
  \item \emph{Recall:} say the Spanish for the lie, then the Spanish for to deceive, then say \emph{honesto} again
\end{itemize}

\begin{tcolorbox}[breakable,colback=teal!4,colframe=teal!35!black,title={Before you move on}]
What does honesto mean? (\textquotedblleft{}Honest\textquotedblright{}.) Say it once more.
\end{tcolorbox}
\section[merecer]{merecer — to deserve}

\label{lesson:ES-C636-merecer}



\begin{quote}
Before the new one: say the Spanish for to deceive, then the Spanish for honest.
\end{quote}

\subsection*{You'll want to know: merecer}
\textbf{merecer} — \textquotedblleft{}to deserve\textquotedblright{}.

\begin{quote}
Mereces unas vacaciones.
\end{quote}

\begin{quote}
\emph{You deserve a holiday.}
\end{quote}

Like \emph{conocer}: \emph{merezco}. \emph{Merecer la pena} — to be worth it.

\begin{cousinweb}
\emph{merecer} is inherited from spoken Latin \emph{merescere}, built on \emph{merere} 'to earn, to deserve' (EtymDB).
\end{cousinweb}

\begin{grammarlens}[title={What you've built}]
One more word for telling, explaining and describing at B1.
\end{grammarlens}

\subsection*{Your turn}
\begin{itemize}
\raggedright
  \item \emph{Say it:} \emph{merecer}
  \item \emph{Say it:} \emph{merecer}, once more
  \item \emph{Say it:} a sentence of your own with \emph{merecer}
  \item \emph{Recall:} say the Spanish for to deceive, then the Spanish for honest, then say \emph{merecer} again
\end{itemize}

\begin{tcolorbox}[breakable,colback=teal!4,colframe=teal!35!black,title={Before you move on}]
What does merecer mean? (\textquotedblleft{}To deserve\textquotedblright{}.) Say it once more.
\end{tcolorbox}
